\section{Data and Experimental Protocol}
\label{sec:data}

\subsection{Corpora}

The task uses two Common Voice \cite{ardila2020commonvoice} releases for Puno
Quechua \cite{huaman2026quechua}. \emph{Scripted Speech 27.0} holds 22{,}727 validated clips, 30.49 hours from 77
speakers, drawn from only 2{,}067 distinct sentences -- so each sentence is read
by roughly eleven different people. That ratio is ordinary for read-speech
collection and it shapes every evaluation decision below.

\emph{Spontaneous Speech 5.0} holds 7{,}305 clips, 35.41 hours from 110
speakers, recorded as free answers to 150 questions. Only 1{,}222 clips carry a
human transcription, 5.81 hours; the remaining 6{,}083 have none, because
transcription requires a Quechua literacy few speakers have. The baseline
release supplies automatic ``silver'' transcriptions for these, generated with
omniASR \cite{keren2025omnilingual}, and we reuse 5{,}782 of them, adding about 21 hours of otherwise unusable speech.

\subsection{The curated reference}
\label{sec:curated}

Written Puno Quechua has no standardised orthography, and contributors spelled
as they saw fit. Scoring against that uneven text would reward reproducing
their inconsistencies rather than hearing well, so the organisers released a
curated reference for the 2{,}067 scripted sentences, and it is what the shared
task scores against. The curation is substantial: 629 sentences differ in spelling beyond case and punctuation,
adopting Puno forms --
\texttt{allqu}$\rightarrow$\texttt{alqu},
\texttt{iskay}$\rightarrow$\texttt{ishkay},
\texttt{misk'i}$\rightarrow$\texttt{mishk'i}. Each is a single-letter change,
and each costs a whole word under WER.

The published baseline predates the curated release and was trained on the raw
text, which isolates the effect exactly: one fixed set of its hypotheses scores
8.13 WER and 2.76 CER against the raw reference, and 16.08 and 3.77 against the
curated one. Nearly eight points of word error rate are orthography rather than
acoustics, and aligning the training text with the curated reference is the
single largest improvement available in this task.

\subsection{Text normalisation}

The curated scripted text uses 27 symbols; the spontaneous transcriptions are
not uniform, with 243 rows carrying 28 characters outside that convention --
digits, braces, brackets, Spanish accents, three Unicode spellings of the
apostrophe and one control character. The apostrophe matters most: in Puno
Quechua it forms part of the letter for ejective consonants (\texttt{k'},
\texttt{p'}, \texttt{ch'}), so three spellings split one phoneme across three
output symbols.

We fold everything to the curated convention: lowercase, Unicode NFC,
\texttt{a}--\texttt{z}, \texttt{ñ} and the apostrophe. Two choices were
verified against the curated reference rather than assumed -- it contains no
accented vowels and no hyphens in 2{,}067 sentences -- leaving a 29-symbol
training alphabet, all of it present in the released checkpoint's 50-way CTC
head.

The folding is orthographic, not lexical: the organisers' curation also
replaces whole word forms (\texttt{allqu}$\rightarrow$\texttt{alqu}), and we
do not propagate those to the spontaneous transcriptions, which nobody curated.
The 372 substitutions recoverable from the scripted diff would touch 1.20\% of
spontaneous training tokens; we leave this untested.

\subsection{The test set does not reuse training sentences}
\label{sec:newsentences}

With only 2{,}067 sentences in the corpus and all of them recorded, it seemed
near certain that the scripted test would reuse them. It does not. Using only
our own greedy hypotheses on the test audio -- no labels -- we measured the
normalised edit distance from each to the nearest sentence in the pool. A
development hypothesis whose sentence is in the pool has a median distance of
0.000; with that sentence removed, 0.516. The test scores 0.480, and its
hypotheses run eight to twelve words against a pool mean of four.

That decides how the rest of this paper is measured: a development set whose
sentences the model has heard from eleven other speakers will prefer systems
that remember text over systems that recognise speech.

\subsection{Splits}
\label{sec:splits}

The baseline's split, which we reproduce as \texttt{sharedtask}, assigns whole
speakers to each partition, balancing gender and duration. It separates
speakers cleanly but cannot separate text: the training portion holds all
2{,}067 sentences, so every development sentence appears there, read by other
people. On the baseline's own released split, 805 of 805 development sentences
and 757 of 757 test sentences also occur in training. This follows from the
corpus rather than from oversight -- using all 22{,}727 clips requires all
2{,}067 sentences in training.

We therefore built a second split, \texttt{sharedtask\_both}, disjoint in
speakers \emph{and} text. Holding out sentences costs clips, so its development
portion falls to 339 scripted clips against 1{,}493; with 1{,}347 reference
words, one WER point is thirteen words, and we treat differences below a few
tenths as noise throughout.

We use both: decisions are made on \texttt{sharedtask\_both}, which matches the
condition the test is in, while the comparison against the published baseline
is reported on \texttt{sharedtask}, the ruler the baseline was measured with.
The spontaneous side needs no reconstruction -- its release carries an official
\texttt{split} column, which we adopt unchanged, excluding the 296 clips marked
as test from training and development alike.

\subsection{The held-out set}
\label{sec:heldout}

A development set that has driven every decision -- ours drove step counts,
learning rates, 72 language model configurations and the decoder per domain --
is no longer an unbiased estimate. None of that is training in the gradient
sense, but all of it is selection, and selection leaves a mark.

We therefore reserved a held-out portion when the splits were built, never
looked at it, and measured it once, after every decision here was frozen. To
make the commitment checkable rather than merely stated, the expected figures
were pre-registered in the evaluation script's header before it ran. Nothing
was changed afterwards.
